\documentclass[11pt]{article}
% Preambulo comun para los informes de la ronda 3.
% Se incluye con % Preambulo comun para los informes de la ronda 3.
% Se incluye con % Preambulo comun para los informes de la ronda 3.
% Se incluye con \input{../comun/preambulo} desde cada subcarpeta.
\usepackage[utf8]{inputenc}
\usepackage[T1]{fontenc}
\usepackage[spanish,es-nodecimaldot]{babel}
\usepackage{lmodern}
\usepackage[a4paper,margin=2.4cm]{geometry}
\usepackage{booktabs}
\usepackage{longtable}
\usepackage{array}
\usepackage{xcolor}
\usepackage{hyperref}
\usepackage{enumitem}
\usepackage{fancyhdr}
\usepackage{titlesec}
\usepackage{graphicx}
\usepackage{amssymb}
\usepackage{microtype}

\definecolor{sosTeal}{HTML}{0F766E}
\definecolor{sosGray}{HTML}{475569}

\hypersetup{
  colorlinks=true,
  linkcolor=sosTeal,
  urlcolor=sosTeal,
  citecolor=sosTeal,
  pdftitle={SOS Clinic - Ronda 3 de pruebas},
  pdfauthor={Equipo SOS Clinic}
}

\pagestyle{fancy}
\fancyhf{}
\lhead{\textcolor{sosGray}{SOS Clinic --- Ronda 3}}
\rhead{\textcolor{sosGray}{\leftmark}}
\cfoot{\thepage}
\setlength{\headheight}{15pt}
\renewcommand{\headrulewidth}{0.4pt}

\titleformat{\section}{\Large\bfseries\color{sosTeal}}{\thesection}{0.8em}{}
\titleformat{\subsection}{\large\bfseries\color{sosGray}}{\thesubsection}{0.6em}{}

\setlist[itemize]{leftmargin=1.4em,itemsep=2pt,topsep=4pt}
\setlist[enumerate]{leftmargin=1.6em,itemsep=2pt,topsep=4pt}

\newcommand{\pantalla}[1]{\texttt{\detokenize{#1}}}
\newcommand{\ok}{\textcolor{sosTeal}{\textbf{PASS}}}
\newcommand{\fallo}{\textcolor{red!70!black}{\textbf{FAIL}}}
\newcommand{\aviso}{\textcolor{orange!80!black}{\textbf{WARN}}}
\newcommand{\omitido}{\textcolor{sosGray}{\textbf{SKIP}}}
 desde cada subcarpeta.
\usepackage[utf8]{inputenc}
\usepackage[T1]{fontenc}
\usepackage[spanish,es-nodecimaldot]{babel}
\usepackage{lmodern}
\usepackage[a4paper,margin=2.4cm]{geometry}
\usepackage{booktabs}
\usepackage{longtable}
\usepackage{array}
\usepackage{xcolor}
\usepackage{hyperref}
\usepackage{enumitem}
\usepackage{fancyhdr}
\usepackage{titlesec}
\usepackage{graphicx}
\usepackage{amssymb}
\usepackage{microtype}

\definecolor{sosTeal}{HTML}{0F766E}
\definecolor{sosGray}{HTML}{475569}

\hypersetup{
  colorlinks=true,
  linkcolor=sosTeal,
  urlcolor=sosTeal,
  citecolor=sosTeal,
  pdftitle={SOS Clinic - Ronda 3 de pruebas},
  pdfauthor={Equipo SOS Clinic}
}

\pagestyle{fancy}
\fancyhf{}
\lhead{\textcolor{sosGray}{SOS Clinic --- Ronda 3}}
\rhead{\textcolor{sosGray}{\leftmark}}
\cfoot{\thepage}
\setlength{\headheight}{15pt}
\renewcommand{\headrulewidth}{0.4pt}

\titleformat{\section}{\Large\bfseries\color{sosTeal}}{\thesection}{0.8em}{}
\titleformat{\subsection}{\large\bfseries\color{sosGray}}{\thesubsection}{0.6em}{}

\setlist[itemize]{leftmargin=1.4em,itemsep=2pt,topsep=4pt}
\setlist[enumerate]{leftmargin=1.6em,itemsep=2pt,topsep=4pt}

\newcommand{\pantalla}[1]{\texttt{\detokenize{#1}}}
\newcommand{\ok}{\textcolor{sosTeal}{\textbf{PASS}}}
\newcommand{\fallo}{\textcolor{red!70!black}{\textbf{FAIL}}}
\newcommand{\aviso}{\textcolor{orange!80!black}{\textbf{WARN}}}
\newcommand{\omitido}{\textcolor{sosGray}{\textbf{SKIP}}}
 desde cada subcarpeta.
\usepackage[utf8]{inputenc}
\usepackage[T1]{fontenc}
\usepackage[spanish,es-nodecimaldot]{babel}
\usepackage{lmodern}
\usepackage[a4paper,margin=2.4cm]{geometry}
\usepackage{booktabs}
\usepackage{longtable}
\usepackage{array}
\usepackage{xcolor}
\usepackage{hyperref}
\usepackage{enumitem}
\usepackage{fancyhdr}
\usepackage{titlesec}
\usepackage{graphicx}
\usepackage{amssymb}
\usepackage{microtype}

\definecolor{sosTeal}{HTML}{0F766E}
\definecolor{sosGray}{HTML}{475569}

\hypersetup{
  colorlinks=true,
  linkcolor=sosTeal,
  urlcolor=sosTeal,
  citecolor=sosTeal,
  pdftitle={SOS Clinic - Ronda 3 de pruebas},
  pdfauthor={Equipo SOS Clinic}
}

\pagestyle{fancy}
\fancyhf{}
\lhead{\textcolor{sosGray}{SOS Clinic --- Ronda 3}}
\rhead{\textcolor{sosGray}{\leftmark}}
\cfoot{\thepage}
\setlength{\headheight}{15pt}
\renewcommand{\headrulewidth}{0.4pt}

\titleformat{\section}{\Large\bfseries\color{sosTeal}}{\thesection}{0.8em}{}
\titleformat{\subsection}{\large\bfseries\color{sosGray}}{\thesubsection}{0.6em}{}

\setlist[itemize]{leftmargin=1.4em,itemsep=2pt,topsep=4pt}
\setlist[enumerate]{leftmargin=1.6em,itemsep=2pt,topsep=4pt}

\newcommand{\pantalla}[1]{\texttt{\detokenize{#1}}}
\newcommand{\ok}{\textcolor{sosTeal}{\textbf{PASS}}}
\newcommand{\fallo}{\textcolor{red!70!black}{\textbf{FAIL}}}
\newcommand{\aviso}{\textcolor{orange!80!black}{\textbf{WARN}}}
\newcommand{\omitido}{\textcolor{sosGray}{\textbf{SKIP}}}


\title{\textbf{DAST con OWASP ZAP}\\[0.3em]
\large Resultados --- Ronda 3}
\author{Equipo SOS Clinic}
\date{22 de octubre de 2026}

\begin{document}
\maketitle

\section{Resumen}

\begin{table}[h]
\centering
\small
\begin{tabular}{@{}lc@{}}
\toprule
Métrica & Valor \\
\midrule
URLs analizadas & 17 \\
Reglas pasadas & \textbf{62 PASS} \\
Fallos & \textbf{0 FAIL} \\
Advertencias & \textbf{6 WARN} \\
Información & 0 INFO \\
\bottomrule
\end{tabular}
\caption{Resumen del baseline de ZAP.}
\end{table}

No hay fallos. Las seis advertencias se evaluaron manualmente; ninguna
representa una vulnerabilidad explotable en el contexto actual.

\section{Advertencias y evaluación}

\begin{table}[h]
\centering
\small
\begin{tabular}{@{}p{4.6cm}cp{7.6cm}@{}}
\toprule
Regla & Ocurrencias & Evaluación \\
\midrule
Cross-Domain JavaScript Source Inclusion [10017] & 3 & Bajo. Recursos de YouTube esperados en el landing; no los controla el usuario. \\
Non-Storable Content [10049] & 9 & Informativo. \texttt{Cache-Control: no-store} es intencional en páginas con sesión. \\
CSP: Wildcard Directive [10055] & 12 & Bajo/medio. Fuentes de esquema y \texttt{unsafe-inline}; se relaciona con F-03. \\
Modern Web Application [10109] & 3 & Informativo. Detecta la SPA; no es vulnerabilidad. \\
Sub Resource Integrity Missing [90003] & 3 & Bajo. Recursos locales del mismo origen; el script de YouTube es externo. \\
Cross-Origin-Embedder-Policy Missing [90004] & 9 & Bajo. COEP es opcional; COOP y CORP sí están presentes. \\
\bottomrule
\end{tabular}
\caption{Advertencias de ZAP y su clasificación.}
\end{table}

\section{Controles que pasaron}

ZAP no encontró problemas en:

\begin{itemize}
  \item Cabeceras anti-clickjacking y \texttt{X-Content-Type-Options}.
  \item Ausencia de CSP (la CSP existe) y divulgación de errores.
  \item Comentarios sospechosos, información sensible en URL o referrer.
  \item Cookies (la aplicación usa \texttt{sessionStorage}, no cookies).
  \item Enlaces a dominios maliciosos, \textit{mixed content} y
  \textit{source code disclosure}.
  \item Tecnologías con versiones vulnerables.
\end{itemize}

Las cabeceras observadas en el frontend incluyen
\texttt{X-Content-Type-Options: nosniff}, \texttt{X-Frame-Options: DENY},
\texttt{Referrer-Policy: no-referrer}, \texttt{Permissions-Policy},
\texttt{COOP}, \texttt{CORP} y una CSP con \texttt{unsafe-inline}.

\section{Mejoras}

\begin{enumerate}
  \item Endurecer la CSP: retirar \texttt{unsafe-inline} y restringir
  \texttt{img-src}/\texttt{media-src} a orígenes concretos.
  \item Añadir \textit{Subresource Integrity} si se incorporan scripts de
  terceros que lo permitan.
  \item Evaluar \texttt{Cross-Origin-Embedder-Policy} si se aíslan recursos
  cross-origin en el futuro.
  \item Mantener el análisis baseline en la CI para detectar regresiones de
  cabeceras.
\end{enumerate}

\section{Evidencias}

\begin{itemize}
  \item \pantalla{seguridad/resultados/zap-baseline-r3.html}
  \item \pantalla{seguridad/resultados/zap-baseline-r3.json}
  \item \pantalla{seguridad/resultados/zap-baseline-r3.md}
\end{itemize}

\section{Conclusión}

El frontend no presenta fallos pasivos y las advertencias son de bajo impacto o
decisiones intencionales. El único punto a mejorar con impacto real es la CSP
con \texttt{unsafe-inline}, ya identificado en el pentest como F-03.

\end{document}
